The canonical unit points and transformed parameters are the arrays
\texttt{sobol\_vectors} and \texttt{parameters} in
\texttt{data/allocation/design.npz}. The generating environment records
Python~3.13.5, NumPy~2.3.5 and SciPy~1.17.0. The exact call is
\begin{verbatim}
from scipy.stats import qmc
u = qmc.Sobol(d=14, scramble=True,
              seed=20260905).random_base2(m=10)
\end{verbatim}
The default \texttt{bits=None} resolves to 30 bits; optimization is
\texttt{None}. No points are skipped or thinned. For the ordered endpoint
vectors in \cref{tab:ranges}, the transformation is
\begin{verbatim}
lo = np.log(lower_endpoints)
hi = np.log(upper_endpoints)
parameters = np.exp(lo + u * (hi - lo))
\end{verbatim}
Every scenario uses the same transformed rows.

Re-execution in SciPy~1.17.0 reproduces the archived unit points bit for
bit. Floating-point exponentiation changes some transformed entries by
at most a relative $2.22\times10^{-16}$ in the present environment;
all classifications agree. Evaluation of the frozen arrays reproduces
219 instantaneous-reference inward corners and 527 under perfect
substitution with zero residual weights, as well as all other counts
in \cref{tab:scenarios}.

With \texttt{seed=20260905}, the archived unit points are reproduced
bit for bit under SciPy~1.17.0 and 1.18.1;
\texttt{rng=20260905} yields a different scramble, with 208
instantaneous and 530 perfect-substitution inward corners.
The frozen arrays, not the seed alone, define the reference experiment.
Supplementary Material section~S2 records the generator specification
and file checksums.

